\documentclass[UTF8]{ctexart}
\usepackage{xeCJK}
\usepackage{geometry}
\geometry{left=2cm,right=2cm,top=2.5cm,bottom=2.5cm}
\setlength{\headheight}{12.7pt}

\usepackage{titlesec}
\titleformat{\section}
  {\normalfont\large\bfseries}
  {\thesection}
  {1em}
  {}
\titlespacing*{\section}
  {0pt}
  {3.5ex plus 1ex minus .2ex}
  {2.3ex plus .2ex}

\usepackage{amsmath}
\usepackage{amssymb}
\usepackage{amsthm}
\usepackage{enumitem}
\setlist[enumerate]{itemsep=0pt,topsep=0pt,parsep=0pt,leftmargin=0pt,itemindent=2.5em}
\setlist[itemize]{itemsep=0pt,topsep=0pt,parsep=0pt,leftmargin=0pt,itemindent=2.5em}
\usepackage{fancyhdr}
\pagestyle{fancy}
\fancyhf{}
\usepackage{graphicx}
\usepackage{hyperref}
\usepackage{indentfirst}
\usepackage{lmodern}


\begin{document}
\setlength{\abovedisplayskip}{1pt}
\setlength{\belowdisplayskip}{1pt}

\section{特征子群}

\hypertarget{sec:定义1.1}{}
\noindent\textbf{定义1.1 (特征子群)}

给定群$G$和它的子群$H$, 如果
\vspace{-3pt}
\[
    \forall\varphi\in\mathrm{Aut}(G):\quad\varphi[H]=H,\\[-3pt]
\]
就称$H$是$G$的\textbf{特征子群}, 记作$H\,\mathrm{char}\,G$.

\vspace{10pt}

\hypertarget{sec:定理1.1}{}
\noindent\textbf{定理1.1}

如果$K\,\mathrm{char}\,H$, $H\,\mathrm{char}\,G$, 那么$K\,\mathrm{char}\,G$.

\noindent{\kaishu 证明.}

对任意的$\varphi\in\mathrm{Aut}(G)$, 有$\varphi[H]=H$,
显然$\varphi\!\restriction_{H}\,\in\mathrm{Aut}(H)$, 从而
$\varphi[K]=\varphi\!\restriction_{H}[K]=K$. \hfill $\qedsymbol$

\vspace{10pt}

\hypertarget{sec:定理1.2}{}
\noindent\textbf{定理1.2}

任给$H<G$, 如果对任意的$\varphi\in\mathrm{Aut}(G)$和$h\in H$
都有$\varphi(h)\in H$, 那么$H\,\mathrm{char}\,G$.

\noindent{\kaishu 证明.}

对任意的$\varphi\in\mathrm{Aut}(G)$ :
\begin{enumerate}
    \item 依假设即得$\varphi[H]\subset H$,
    \item 因为$\varphi^{-1}$也是自同构, 所以$\varphi^{-1}[H]\subset H$,
    从而$H\subset\varphi[H]$,
\end{enumerate}
综上$\varphi[H]=H$. \hfill $\qedsymbol$

\vspace{15pt}

下面举一些特征子群的例子.

\vspace{10pt}

\hypertarget{sec:例1.1}{}
\noindent\textbf{例1.1}

任给群$G$, 它的中心$Z(G)$是$G$的特征子群.

\noindent{\kaishu 证明.}

对任意的$\varphi\in\mathrm{Aut}(G)$, $g\in Z(G)$和$h\in G$, 有
\[
    \varphi(g)\cdot h=\varphi(g\cdot\varphi^{-1}h)=
    \varphi(\varphi^{-1}h\cdot g)=h\cdot\varphi(g),
\]
从而$\varphi(g)\in Z(G)$. 所以$Z(G)\,\mathrm{char}\,G$. \hfill $\qedsymbol$

\vspace{10pt}

\hypertarget{sec:例1.2}{}
\noindent\textbf{例1.2}

任给群$G$, 对任意的$g,h\in G$, 定义
\vspace{-3pt}
\[
    [g,h]\overset{\mathrm{def}}{=}g^{-1}h^{-1}gh,\\[-3pt]
\]
再定义$G'\overset{\mathrm{def}}{=}\langle\,\{[g,h]\mid g,h\in G\}\,\rangle$,
称为$G$的\textbf{换位子群}. 它是$G$的特征子群.

\noindent{\kaishu 证明.}

首先, 显然对任意的$\varphi\in\mathrm{Aut}(G)$和$g,h\in G$有
\[
    \varphi[g,h]=[\varphi(g),\varphi(h)],\quad[g,h]^{-1}=[h,g],
\]
从而对任意的$z\in G'$, 存在自然数$n$和$x,y:n\to G$使得
\vspace{3pt}
\[
    z=\prod_{i=0}^n\,[x_i,y_i]\quad\implies\quad
    \varphi(z)=\prod_{i=0}^n\,[\varphi(x_i),\varphi(y_i)]\in G',
\]
所以$G'\,\mathrm{char}\,G$. \hfill $\qedsymbol$

\end{document}